\documentclass{zhvt-classic}

% 使用本地源流明体（GenRyuMin2 TW）按文件路径加载，不依赖系统字体名
% 版框（内框／界栏／书口侧栏／鱼尾）改用与版心印章同一印泥棕色：
%   印章图片 yshs.png 实色 = RGB(135,68,52) = HTML 874434
\usepackage{xcolor}
\usepackage{tikz}
% sealbrown 双注册：xcolor 供 TikZ/pgf，L3 供 cls 版框
\definecolor{sealbrown}{HTML}{874434}
\ExplSyntaxOn
\color_set:nnn { sealbrown } { HTML } { 874434 }
\ExplSyntaxOff

\zhvtset{
  grid color       = sealbrown,
  main font        = GenRyuMin2TW-SB,
  main font path   = C:/Users/Administrator/texmf/fonts/opentype/public/GenRyuMin2TW/,
jiazhu font      = GenRyuMin2TW-M,
  jiazhu font path = C:/Users/Administrator/texmf/fonts/opentype/public/GenRyuMin2TW/,
  % 页码框容纳 4 字：正常页「九十七」+前置空 1 格=4 行恰好；
  % 100+ 页「一百零一」共 5 行超出，末字「一/二/三」被框裁掉（2026-10-04 实测），故放宽到 5
  shukou page number chars = 5,
  % 书口书名框高从默认 6 字减到 4 字，把省出的 2 字高拨给 chapter 框，
  % 使 6 字书名+卷名放进 9 字高的 chapter 框后不再压线（2026-10-05）
  shukou title chars = 4,
}

\title{史记}[高祖本纪]
\author{汉·司马迁}
\maker{闲叟漫记·重製}


% ===== 索隐述赞样式（2026-11-08 四次定稿 · 愚大人亲手改定参数）=====
% ⚠ 竖排前提：自定义字体家族必须带 RawFeature={vertical:+vert}；
%   simhei.ttf 加 +vhal 会整体加载失败静默回退，只写 +vert；simfang 同理。
\newCJKfontfamily{\shuzanfont}[Path=C:/Windows/Fonts/,RawFeature={vertical:+vert}]{simhei.ttf}
\newCJKfontfamily{\shuzanbody}[Path=C:/Windows/Fonts/,RawFeature={vertical:+vert}]{simfang.ttf}
% 题头：152x30pt 圆角 6pt 棕底白字；baseline (0,2) 使整框落入界行格子不压版框线；
% 白字节点 x 不补偿、y +8.2pt（像素级墨迹实测 dx/dy≈0，L/R 留白约 20pt）。
\newcommand{\shuzanhead}[1]{%
  \begin{tikzpicture}[baseline={(0,2)}]
    \fill[sealbrown, rounded corners=6pt]
      (-76pt,-15pt) rectangle (76pt,15pt);
    \node[text=white, inner sep=0pt, anchor=center] at (0pt,8.2pt)
      {\shuzanfont\fontsize{23}{23}\selectfont #1};
  \end{tikzpicture}%
}
% ★ 修正书口卷名位置：chapter 框起始位置改回 cls 默认值，避免全局影响。
%   只在每个卷首页通过 \juan 命令里的空命令微调该页位置（安全，不会跨页污染）。
\ExplSyntaxOn
\AtBeginDocument{
  \dim_set:Nn \l__zhvt_shukou_chapter_title_start_dim
    { \l__zhvt_shukou_title_height_dim + \l__zhvt_font_size_dim * 2 }
  \dim_set:Nn \l__zhvt_shukou_page_number_start_dim
    { \l__zhvt_shukou_chapter_title_start_dim + \l__zhvt_shukou_chapter_title_height_dim }
}
\ExplSyntaxOff

\begin{document}

% ★★ 版心（书口）排布覆写：把书名放到鱼尾之下（2026-10-04）
%
%   cls 默认版心自上而下：书名（title 静态框，在鱼尾之上）→ 鱼尾 → 篇名（chapter 动态框）→ 页码
%   本模板改为：          鱼尾（位置不动）→ 书名 → 篇名 → 页码分割线 → 页码
%
%   手法：把 cls 放在 title 静态框里的书名，改放进 chapter 动态框（鱼尾之下）；
%         title 静态框留空。书名取 \l__zhvt_title_str，即 \title{...} 的第一参数，
%         所以换书时只改 \title，不必在此重复写书名。
%
%   ⚠ 三条必须遵守的写法（2026-10-04 实测）：
%     ① 汉字必须用 \hbox{...} 包住。flowfram 会把框内容装进一个宽 14.4mm（= 一个 line_sep
%        = 书口栏宽）的纵向 minipage 并按该宽度折行，而整页靠 PDF /Rotate 90 显示、
%        框内文字本身不变向 —— 不包 \hbox 时「围炉夜话」4 字会被逐字折成 4 行，
%        在版面上就成了横排 4 字并横向溢出版框。
%     ② 竖排下**一个 \hbox 就是一整列**。书名与篇名要同处一列，必须放进同一个 \hbox；
%        分成两个 \hbox 时，后一个会被当成新的一列、甩到隔壁栏去
%        （实测「序」跑到书口栏外侧 14.4mm 的正文列上）。中间的 \quad 在框内是横向空白，
%        显示出来正好是书名与篇名之间的纵向间隔。
%     ③ ★ 内容盒要写成 **与框等宽的定宽盒** `\hbox to \l__zhvt_shukou_chapter_title_height_dim{…\hfil}`。
%        实测：chapter 框是把内容盒沿纵向**居中**摆放的，内容宽度一变，书名就跟着漂：
%        自然宽度的 \hbox（书名 38.4mm）落在 x=157.2mm，加篇名后（57.6mm）落在 x=152.4mm，
%        目录页（无篇名）与正文页的书名就差半格。做成等宽盒后，无论有没有篇名，
%        内容固定从框顶起排，实测两页均落在 x=142.8mm，不再漂。
\ExplSyntaxOn
% 把书口 chapter 框（书名+篇名）起始位置下移 1 字高，避免 4 字卷名顶到鱼尾
% ⚠ cls 的 \__zhvt_calculate_dimensions: 在 \AtBeginDocument 里执行，会重算该 dim；
%   必须把覆写挂到 begin document 之后，否则会被二次覆盖（2026-10-05 实测）。
% (旧方案已移到导言区，直接修改 chapter 框起始位置)
% 卷标题命令 #1=书口篇名(鱼尾) #2=正文卷名（含卷序，卷名内空格用 ~）
%
% ★★ 目录页两项修订（2026-10-05，愚大人指出）★★
%   ① 页码起点统一对齐：
%      原状 —— 目录项写成 \par\hspace{4em}卷名\hspace{4em}页码，页码起点 = 卷名实宽 + 4em。
%      卷名实宽随字数变：「卷一~梁惠王上」6 字 +~、「卷七~离娄上」5 字 +~，
%      差 1 个字位 → 卷七~卷十 4 条的页码整体比其余 10 条高 27.21pt（一个字），看着上下不齐。
%      改法 —— 卷名补全角空格（宽 = 1 字 = 27.215pt）补足到 **7 字符位**再写进目录，
%      14 条页码起点遂全部落在同一条竖线上（起点对齐；末字随页码字数自然参差）。
%      ⚠ 只对目录里的副本 \l_tmpc_tl 补，正文卷名仍用 \l_tmpa_tl 原样 —— 正文像素不受影响。
%   ② 卷名整体上移半字位：前导 \hspace 由 4em 改 3.5em。
%      实测原状目录块在版框内偏下（天头留白 139.6pt / 地脚留白 110.7pt），
%      上移 0.5em = 13.6pt 后上下基本均分，视觉居中。
\NewDocumentCommand { \juan } { mm } {
	\clearpage
	\tl_set:Nn \l__zhvt_chapter_title_tl { #1 }
	\tl_set:Nn \l_tmpa_tl { #2 }
	\tl_set_eq:NN \l_tmpc_tl \l_tmpa_tl
	\int_step_inline:nn { 7 - \str_count:n { #2 } }
	  { \tl_put_right:Nn \l_tmpc_tl { 　 } }
	\tl_set:Nx \l_tmpb_tl { \zhnum { page } }
	\tl_put_left:Nn \l_tmpb_tl { \hspace { 4em } }
	\tl_put_left:Nn \l_tmpb_tl { \l_tmpc_tl }
	\tl_put_left:Nn \l_tmpb_tl { \par \hspace { 3.5em } }
	\addtocontents { toc } { \l_tmpb_tl }
	\skip_horizontal:n { \l__zhvt_font_size_dim * 4 }
	\l_tmpa_tl
	\par
	\thispagestyle { empty }
}
\cs_set_protected:Npn \__zhvt_set_shukou_content: {
	% ★ 鱼尾上方留白（title 静态框）嵌入闲叟钤印：
	%   整页 /Rotate 90，故此处把图用 \rotatebox{90} 转正为竖章；
	%   尺寸卡在版心条幅内（书内约 7mm 宽 × 33.8mm 高，留余量不压鱼尾）。
	\setstaticcontents * { title } {
		\vbox_to_ht:nn { \l__zhvt_baseline_skip_dim } {
			\vskip 14.2pt
			\hbox_to_wd:nn { \l__zhvt_shukou_title_height_dim } {
				\hfil
				\rotatebox [ origin = c ] { 90 } { \includegraphics [ width = 7mm , height = 33.8mm ] { yshs.png } }
				\hfil
			}
		}
	}
% ★ 版心（书口）统一改法（2026-11-16 愚大人定）：18pt＋框内居中＋书口文字同版框色
%   ① 原 27.2pt 时 7 字长篇名（如「廉颇蔺相如列传」）末字溢出 chapter 框、压住页码；
%      改 \fontsize{18pt}{18pt}，前后 \hfil 令字串在框内居中（不随篇名字数漂移）。
%   ② \vbox{\vskip 8.2pt} 定量下移，抵消缩字号后栏宽方向 8.2pt 上偏（像素实测 y 中心 543.3 = 原版）。
%      ⚠ \raisebox / \vspace* 在 flowfram 框内无效；换字号须同步换补偿：16pt→10.1、19pt→7.4。
%   ③ chapter 与 page~number 均置 \color{sealbrown}（HTML 874434），与版框／界行同色；正文仍黑。
%   ⚠ 页码字位、钤印、鱼尾、界行、正文一律未动；五部合订本同此口径。
	\setdynamiccontents * { chapter }
	{ \color{sealbrown} \vbox { \vskip 8.2pt \hbox to \l__zhvt_shukou_chapter_title_height_dim
		{ \hfil \fontsize{18pt}{18pt}\selectfont \l__zhvt_title_str \quad \l__zhvt_chapter_title_tl \hfil } } }
	\setdynamiccontents * { page~number } { \color{sealbrown} \quad \zhnum { page } }
}
\__zhvt_set_shukou_content:
\ExplSyntaxOff

\maketitle

% 单篇，无多卷目录，不渲染目录页




% 方案①：仅改 A 类错讹，头部与版式不变

\juan{高祖本纪}{卷八~高祖本纪第八}
\hspace*{2em}高祖。沛丰邑中阳里人。姓刘氏。字季。父曰太公。母曰刘媪。其先刘媪尝息大泽之陂。梦与神遇。是时雷电晦冥。太公往视。则见蛟龙於其上。已而有身。遂产高祖。\\
\hspace*{2em}高祖为人。隆准而龙颜。美须髯。左股有七十二黑子。仁而爱人。喜施。意豁如也。常有大度。不事家人生产作业。及壮。试为吏。为泗水亭长。廷中吏无所不狎侮。好酒及色。常从王媪。武负贳酒。醉卧。武负。王媪见其上常有龙。怪之。高祖每酤留饮。酒雠数倍。及见怪。岁竟。此两家常折券弃责。\\
\hspace*{2em}高祖常繇咸阳。纵观。观秦皇帝。喟然太息曰。﹁嗟乎。大丈夫当如此也。﹂\\
\hspace*{2em}单父人吕公善沛令。避仇从之客。因家沛焉。沛中豪桀吏闻令有重客。皆往贺。萧何为主吏。主进。令诸大夫曰。﹁进不满千钱。坐之堂下。﹂高祖为亭长。素易诸吏。乃绐为谒曰﹁贺钱万﹂。实不持一钱。谒入。吕公大惊。起。迎之门。吕公者。好相人。见高祖状貌。因重敬之。引入坐。萧何曰。﹁刘季固多大言。少成事。﹂高祖因狎侮诸客。遂坐上坐。无所诎。酒阑。吕公因目固留高祖。高祖竟酒。後。吕公曰。﹁臣少好相人。相人多矣。无如季相。原季自爱。臣有息女。原为季箕帚妾。﹂酒罢。吕媪怒吕公曰。﹁公始常欲奇此女。与贵人。沛令善公。求之不与。何自妄许与刘季。﹂吕公曰。﹁此非兒女子所知也。﹂卒与刘季。吕公女乃吕后也。生孝惠帝。鲁元公主。\\
\hspace*{2em}高祖为亭长时。常告归之田。吕后与两子居田中耨。有一老父过请饮。吕后因餔之。老父相吕后曰。﹁夫人天下贵人。﹂令相两子。见孝惠。曰。﹁夫人所以贵者。乃此男也。﹂相鲁元。亦皆贵。老父已去。高祖適从旁舍来。吕后具言客有过。相我子母皆大贵。高祖问。曰。﹁未远。﹂乃追及。问老父。老父曰。﹁乡者夫人婴兒皆似君。君相贵不可言。﹂高祖乃谢曰。﹁诚如父言。不敢忘德。﹂及高祖贵。遂不知老父处。\\
\hspace*{2em}高祖为亭长。乃以竹皮为冠。令求盗之薛治之。时时冠之。及贵常冠。所谓﹁刘氏冠﹂乃是也。\\
\hspace*{2em}高祖以亭长为县送徒骊山。徒多道亡。自度比至皆亡之。到丰西泽中。止饮。夜乃解纵所送徒。曰。﹁公等皆去。吾亦从此逝矣。﹂徒中壮士愿从者十馀人。高祖被酒。夜径泽中。令一人行前。行前者还报曰。﹁前有大蛇当径。愿还。﹂高祖醉。曰。﹁壮士行。何畏。﹂乃前。拔剑击斩蛇。蛇遂分为两。径开。行数里。醉。因卧。後人来至蛇所。有一老妪夜哭。人问何哭。妪曰。﹁人杀吾子。故哭之。﹂人曰。﹁妪子何为见杀。﹂妪曰。﹁吾子。白帝子也。化为蛇。当道。今为赤帝子斩之。故哭。﹂人乃以妪为不诚。欲笞之。妪因忽不见。後人至。高祖觉。後人告高祖。高祖乃心独喜。自负。诸从者日益畏之。\\
\hspace*{2em}秦始皇帝常曰﹁东南有天子气﹂。於是因东游以厌之。高祖即自疑。亡匿。隐於芒。砀山泽岩石之间。吕后与人俱求。常得之。高祖怪问之。吕后曰。﹁季所居上常有云气。故从往常得季。﹂高祖心喜。沛中子弟或闻之。多欲附者矣。\\
\hspace*{2em}秦二世元年秋。陈胜等起蕲。至陈而王。号为﹁张楚﹂。诸郡县皆多杀其长吏以应陈涉。沛令恐。欲以沛应涉。掾。主吏萧何。曹参乃曰。﹁君为秦吏。今欲背之。率沛子弟。恐不听。愿君召诸亡在外者。可得数百人。因劫众。众不敢不听。﹂乃令樊哙召刘季。刘季之众已数十百人矣。\\
\hspace*{2em}於是樊哙从刘季来。沛令後悔。恐其有变。乃闭城城守。欲诛萧。曹。萧。曹恐。逾城保刘季。刘季乃书帛射城上。谓沛父老曰。﹁天下苦秦久矣。今父老虽为沛令守。诸侯并起。今屠沛。沛今共诛令。择子弟可立者立之。以应诸侯。则家室完。不然。父子俱屠。无为也。﹂父老乃率子弟共杀沛令。开城门迎刘季。欲以为沛令。刘季曰。﹁天下方扰。诸侯并起。今置将不善。壹败涂地。吾非敢自爱。恐能薄。不能完父兄子弟。此大事。愿更相推择可者。﹂萧。曹等皆文吏。自爱。恐事不就。後秦种族其家。尽让刘季。诸父老皆曰。﹁平生所闻刘季诸珍怪。当贵。且卜筮之。莫如刘季最吉。﹂於是刘季数让。众莫敢为。乃立季为沛公。祠黄帝。祭蚩尤於沛庭。而衅鼓旗。帜皆赤。由所杀蛇白帝子。杀者赤帝子。故上赤。於是少年豪吏如萧。曹。樊哙等皆为收沛子弟二三千人。攻胡陵。方与。还守丰。\\
\hspace*{2em}秦二世二年。陈涉之将周章军西至戏而还。燕。赵。齐。魏皆自立为王。项氏起吴。秦泗州监平将兵围丰。二日。出与战。破之。命雍齿守丰。引兵之薛。泗州守壮败於薛。走至戚。沛公左司马得泗州守壮。杀之。沛公还军亢父。至方与。未战。陈王使魏人周巿略地。周巿使人谓雍齿曰。﹁丰。故梁徙也。今魏地已定者数十城。齿今下魏。魏以齿为侯守丰。不下。且屠丰。﹂雍齿雅不欲属沛公。及魏招之。即反为魏守丰。沛公引兵攻丰。不能取。沛公病。还之沛。沛公怨雍齿与丰子弟叛之。闻东阳甯君。秦嘉立景驹为假王。在留。乃往从之。欲请兵以攻丰。是时秦将章邯从陈。别将司马夷将兵北定楚地。屠相。至砀。东阳甯君。沛公引兵西。与战萧西。不利。还收兵聚留。引兵攻砀。三日乃取砀。因收砀兵。得五六千人。攻下邑。拔之。还军丰。闻项梁在薛。从骑百馀往见之。项梁益沛公卒五千人。五大夫将十人。沛公还。引兵攻丰。\\
\hspace*{2em}从项梁月馀。项羽已拔襄城还。项梁尽召别将居薛。闻陈王定死。因立楚後怀王孙心为楚王。治盱台。项梁号武信君。居数月。北攻亢父。救东阿。破秦军。齐军归。楚独追北。使沛公。项羽别攻城阳。屠之。军濮阳之东。与秦军战。破之。\\
\hspace*{2em}秦军复振。守濮阳。环水。楚军去而攻定陶。定陶未下。沛公与项羽西略地至雍丘之下。与秦军战。大破之。斩李由。还攻外黄。外黄未下。\\
\hspace*{2em}项梁再破秦军。有骄色。宋义谏。不听。秦益章邯兵。夜衔枚击项梁。大破之定陶。项梁死。沛公与项羽方攻陈留。闻项梁死。引兵与吕将军俱东。吕臣军彭城东。项羽军彭城西。沛公军砀。\\
\hspace*{2em}章邯已破项梁军。则以为楚地兵不足忧。乃渡河。北击赵。大破之。当是之时。赵歇为王。秦将王离围之钜鹿城。此所谓河北之军也。\\
\hspace*{2em}秦二世三年。楚怀王见项梁军破。恐。徙盱台都彭城。并吕臣。项羽军自将之。以沛公为砀郡长。封为武安侯。将砀郡兵。封项羽为长安侯。号为鲁公。吕臣为司徒。其父吕青为令尹。\\
\hspace*{2em}赵数请救。怀王乃以宋义为上将军。项羽为次将。范增为末将。北救赵。令沛公西略地入关。与诸将约。先入定关中者王之。\\
\hspace*{2em}当是时。秦兵彊。常乘胜逐北。诸将莫利先入关。独项羽怨秦破项梁军。奋。愿与沛公西入关。怀王诸老将皆曰。﹁项羽为人僄悍猾贼。项羽尝攻襄城。襄城无遗类。皆阬之。诸所过无不残灭。且楚数进取。前陈王。项梁皆败。不如更遣长者扶义而西。告谕秦父兄。秦父兄苦其主久矣。今诚得长者往。毋侵暴。宜可下。今项羽僄悍。今不可遣。独沛公素宽大长者。可遣。﹂卒不许项羽。而遣沛公西略地。收陈王。项梁散卒。乃道砀至成阳。与杠里秦军夹壁。破二军。楚军出兵击王离。大破之。\\
\hspace*{2em}沛公引兵西。遇彭越昌邑。因与俱攻秦军。战不利。还至栗。遇刚武侯。夺其军。可四千馀人。并之。与魏将皇欣。魏申徒武蒲之军并攻昌邑。昌邑未拔。西过高阳。郦食其监门。曰。﹁诸将过此者多。吾视沛公大人长者。﹂乃求见说沛公。沛公方踞床。使两女子洗足。郦生不拜。长揖。曰。﹁足下必欲诛无道秦。不宜踞见长者。﹂於是沛公起。摄衣谢之。延上坐。食其说沛公袭陈留。得秦积粟。乃以郦食其为广野君。郦商为将。将陈留兵。与偕攻开封。开封未拔。西与秦将杨熊战白马。又战曲遇东。大破之。杨熊走之荥阳。二世使使者斩以徇。南攻颍阳。屠之。因张良遂略韩地轘辕。\\
\hspace*{2em}当是时。赵别将司马卬方欲渡河入关。沛公乃北攻平阴。绝河津。南。战雒阳东。军不利。还至阳城。收军中马骑。与南阳守齮战犨东。破之。略南阳郡。南阳守齮走。保城守宛。沛公引兵过而西。张良谏曰。﹁沛公虽欲急入关。秦兵尚众。距险。今不下宛。宛从後击。彊秦在前。此危道也。﹂於是沛公乃夜引兵从他道还。更旗帜。黎明。围宛城三匝。南阳守欲自刭。其舍人陈恢曰。﹁死未晚也。﹂乃逾城见沛公。曰。﹁臣闻足下约。先入咸阳者王之。今足下留守宛。宛。大郡之都也。连城数十。人民众。积蓄多。吏人自以为降必死。故皆坚守乘城。今足下尽日止攻。士死伤者必多。引兵去宛。宛必随足下後。足下前则失咸阳之约。後又有彊宛之患。为足下计。莫若约降。封其守。因使止守。引其甲卒与之西。诸城未下者。闻声争开门而待。足下通行无所累。﹂沛公曰。﹁善。﹂乃以宛守为殷侯。封陈恢千户。引兵西。无不下者。至丹水。高武侯鰓。襄侯王陵降西陵。还攻胡阳。遇番君别将梅鋗。与皆。降析。郦。遣魏人甯昌使秦。使者未来。是时章邯已以军降项羽於赵矣。\\
\hspace*{2em}初。项羽与宋义北救赵。及项羽杀宋义。代为上将军。诸将黥布皆属。破秦将王离军。降章邯。诸侯皆附。及赵高已杀二世。使人来。欲约分王关中。沛公以为诈。乃用张良计。使郦生。陆贾往说秦将。啗以利。因袭攻武关。破之。又与秦军战於蓝田南。益张疑兵旗帜。诸所过毋得掠卤。秦人憙。秦军解。因大破之。又战其北。大破之。乘胜。遂破之。\\
\hspace*{2em}汉元年十月。沛公兵遂先诸侯至霸上。秦王子婴素车白马。系颈以组。封皇帝玺符节。降轵道旁。诸将或言诛秦王。沛公曰。﹁始怀王遣我。固以能宽容。且人已服降。又杀之。不祥。﹂乃以秦王属吏。遂西入咸阳。欲止宫休舍。樊哙。张良谏。乃封秦重宝财物府库。还军霸上。召诸县父老豪桀曰。﹁父老苦秦苛法久矣。诽谤者族。偶语者弃巿。吾与诸侯约。先入关者王之。吾当王关中。与父老约。法三章耳。杀人者死。伤人及盗抵罪。馀悉除去秦法。诸吏人皆案堵如故。凡吾所以来。为父老除害。非有所侵暴。无恐。且吾所以还军霸上。待诸侯至而定约束耳。﹂乃使人与秦吏行县乡邑。告谕之。秦人大喜。争持牛羊酒食献飨军士。沛公又让不受。曰。﹁仓粟多。非乏。不欲费人。﹂人又益喜。唯恐沛公不为秦王。\\
\hspace*{2em}或说沛公曰。﹁秦富十倍天下。地形彊。今闻章邯降项羽。项羽乃号为雍王。王关中。今则来。沛公恐不得有此。可急使兵守函谷关。无内诸侯军。稍徵关中兵以自益。距之。﹂沛公然其计。从之。十一月中。项羽果率诸侯兵西。欲入关。关门闭。闻沛公已定关中。大怒。使黥布等攻破函谷关。十二月中。遂至戏。沛公左司马曹无伤闻项王怒。欲攻沛公。使人言项羽曰。﹁沛公欲王关中。令子婴为相。珍宝尽有之。﹂欲以求封。亚父劝项羽击沛公。方飨士。旦日合战。是时项羽兵四十万。号百万。沛公兵十万。号二十万。力不敌。会项伯欲活张良。夜往见良。因以文谕项羽。项羽乃止。沛公从百馀骑。驱之鸿门。见谢项羽。项羽曰。﹁此沛公左司马曹无伤言之。不然。籍何以至此。﹂沛公以樊哙。张良故。得解归。归。立诛曹无伤。\\
\hspace*{2em}项羽遂西。屠烧咸阳秦宫室。所过无不残破。秦人大失望。然恐。不敢不服耳。\\
\hspace*{2em}项羽使人还报怀王。怀王曰。﹁如约。﹂项羽怨怀王不肯令与沛公俱西入关。而北救赵。後天下约。乃曰。﹁怀王者。吾家项梁所立耳。非有功伐。何以得主约。本定天下。诸将及籍也。﹂乃详尊怀王为义帝。实不用其命。\\
\hspace*{2em}正月。项羽自立为西楚霸王。王梁。楚地九郡。都彭城。负约。更立沛公为汉王。王巴。蜀。汉中。都南郑。三分关中。立秦三将。章邯为雍王。都废丘。司马欣为塞王。都栎阳。董翳为翟王。都高奴。楚将瑕丘申阳为河南王。都洛阳。赵将司马卬为殷王。都朝歌。赵王歇徙王代。赵相张耳为常山王。都襄国。当阳君黥布为九江王。都六。怀王柱国共敖为临江王。都江陵。番君吴芮为衡山王。都邾。燕将臧荼为燕王。都蓟。故燕王韩广徙王辽东。广不听。臧荼攻杀之无终。封成安君陈馀河间三县。居南皮。封梅鋗十万户。\\
\hspace*{2em}四月。兵罢戏下。诸侯各就国。汉王之国。项王使卒三万人从。楚与诸侯之慕从者数万人。从杜南入蚀中。去辄烧绝栈道。以备诸侯盗兵袭之。亦示项羽无东意。至南郑。诸将及士卒多道亡归。士卒皆歌思东归。韩信说汉王曰。﹁项羽王诸将之有功者。而王独居南郑。是迁也。军吏士卒皆山东之人也。日夜跂而望归。及其锋而用之。可以有大功。天下已定。人皆自宁。不可复用。不如决策东乡。争权天下。﹂\\
\hspace*{2em}项羽出关。使人徙义帝。曰。﹁古之帝者地方千里。必居上游。﹂乃使使徙义帝长沙郴县。趣义帝行。群臣稍倍叛之。乃阴令衡山王。临江王击之。杀义帝江南。项羽怨田荣。立齐将田都为齐王。田荣怒。因自立为齐王。杀田都而反楚。予彭越将军印。令反梁地。楚令萧公角击彭越。彭越大破之。陈馀怨项羽之弗王己也。令夏说说田荣。请兵击张耳。齐予陈馀兵。击破常山王张耳。张耳亡归汉。迎赵王歇於代。复立为赵王。赵王因立陈馀为代王。项羽大怒。北击齐。\\
\hspace*{2em}八月。汉王用韩信之计。从故道还。袭雍王章邯。邯迎击汉陈仓。雍兵败。还走。止战好畤。又复败。走废丘。汉王遂定雍地。东至咸阳。引兵围雍王废丘。而遣诸将略定陇西。北地。上郡。令将军薛欧。王西出武关。因王陵兵南阳。以迎太公。吕后於沛。楚闻之。发兵距之阳夏。不得前。令故吴令郑昌为韩王。距汉兵。\\
\hspace*{2em}二年。汉王东略地。塞王欣。翟王翳。河南王申阳皆降。韩王昌不听。使韩信击破之。於是置陇西。北地。上郡。渭南。河上。中地郡。关外置河南郡。更立韩太尉信为韩王。诸将以万人若以一郡降者。封万户。缮治河上塞。诸故秦苑囿园池。皆令人得田之。正月。虏雍王弟章平。大赦罪人。\\
\hspace*{2em}汉王之出关至陕。抚关外父老。还。张耳来见。汉王厚遇之。\\
\hspace*{2em}二月。令除秦社稷。更立汉社稷。\\
\hspace*{2em}三月。汉王从临晋渡。魏王豹将兵从。下河内。虏殷王。置河内郡。南渡平阴津。至雒阳。新城三老董公遮说汉王以义帝死故。汉王闻之。袒而大哭。遂为义帝发丧。临三日。发使者告诸侯曰。﹁天下共立义帝。北面事之。今项羽放杀义帝於江南。大逆无道。寡人亲为发丧。诸侯皆缟素。悉发关内兵。收三河士。南浮江汉以下。愿从诸侯王击楚之杀义帝者。﹂\\
\hspace*{2em}是时项王北击齐。田荣与战城阳。田荣败。走平原。平原民杀之。齐皆降楚。楚因焚烧其城郭。系虏其子女。齐人叛之。田荣弟横立荣子广为齐王。齐王反楚城阳。项羽虽闻汉东。既已连齐兵。欲遂破之而击汉。汉王以故得劫五诸侯兵。遂入彭城。项羽闻之。乃引兵去齐。从鲁出胡陵。至萧。与汉大战彭城灵壁东睢水上。大破汉军。多杀士卒。睢水为之不流。乃取汉王父母妻子於沛。置之军中以为质。当是时。诸侯见楚彊汉败。还皆去汉复为楚。塞王欣亡入楚。\\
\hspace*{2em}吕后兄周吕侯为汉将兵。居下邑。汉王从之。稍收士卒。军砀。汉王乃西过梁地。至虞。使谒者随何之九江王布所。曰。﹁公能令布举兵叛楚。项羽必留击之。得留数月。吾取天下必矣。﹂随何往说九江王布。布果背楚。楚使龙且往击之。\\
\hspace*{2em}汉王之败彭城而西。行使人求家室。家室亦亡。不相得。败後乃独得孝惠。六月。立为太子。大赦罪人。令太子守栎阳。诸侯子在关中者皆集栎阳为卫。引水灌废丘。废丘降。章邯自杀。更名废丘为槐里。於是令祠官祀天地四方上帝山川。以时祀之。兴关内卒乘塞。\\
\hspace*{2em}是时九江王布与龙且战。不胜。与随何间行归汉。汉王稍收士卒。与诸将及关中卒益出。是以兵大振荥阳。破楚京。索间。\\
\hspace*{2em}三年。魏王豹谒归视亲疾。至即绝河津。反为楚。汉王使郦生说豹。豹不听。汉王遣将军韩信击。大破之。虏豹。遂定魏地。置三郡。曰河东。太原。上党。汉王乃令张耳与韩信遂东下井陉击赵。斩陈馀。赵王歇。其明年。立张耳为赵王。\\
\hspace*{2em}汉王军荥阳南。筑甬道属之河。以取敖仓。与项羽相距岁馀。项羽数侵夺汉甬道。汉军乏食。遂围汉王。汉王请和。割荥阳以西者为汉。项王不听。汉王患之。乃用陈平之计。予陈平金四万斤。以间疏楚君臣。於是项羽乃疑亚父。亚父是时劝项羽遂下荥阳。及其见疑。乃怒。辞老。愿赐骸骨归卒伍。未至彭城而死。\\
\hspace*{2em}汉军绝食。乃夜出女子东门二千馀人。被甲。楚因四面击之。将军纪信乃乘王驾。诈为汉王。诳楚。楚皆呼万岁。之城东观。以故汉王得与数十骑出西门遁。令御史大夫周苛。魏豹。枞公守荥阳。诸将卒不能从者。尽在城中。周苛。枞公相谓曰。﹁反国之王。难与守城。﹂因杀魏豹。\\
\hspace*{2em}汉王之出荥阳入关。收兵欲复东。袁生说汉王曰。﹁汉与楚相距荥阳数岁。汉常困。愿君王出武关。项羽必引兵南走。王深壁。令荥阳成皋间且得休。使韩信等辑河北赵地。连燕齐。君王乃复走荥阳。未晚也。如此。则楚所备者多。力分。汉得休。复与之战。破楚必矣。﹂汉王从其计。出军宛叶间。与黥布行收兵。\\
\hspace*{2em}项羽闻汉王在宛。果引兵南。汉王坚壁不与战。是时彭越渡睢水。与项声。薛公战下邳。彭越大破楚军。项羽乃引兵东击彭越。汉王亦引兵北军成皋。项羽已破走彭越。闻汉王复军成皋。乃复引兵西。拔荥阳。诛周苛。枞公。而虏韩王信。遂围成皋。\\
\hspace*{2em}汉王跳。独与滕公共车出成皋玉门。北渡河。驰宿脩武。自称使者。晨驰入张耳。韩信壁。而夺之军。乃使张耳北益收兵赵地。使韩信东击齐。汉王得韩信军。则复振。引兵临河。南飨军小脩武南。欲复战。郎中郑忠乃说止汉王。使高垒深堑。勿与战。汉王听其计。使卢绾。刘贾将卒二万人。骑数百。渡白马津。入楚地。与彭越复击破楚军燕郭西。遂复下梁地十馀城。\\
\hspace*{2em}淮阴已受命东。未渡平原。汉王使郦生往说齐王田广。广叛楚。与汉和。共击项羽。韩信用蒯通计。遂袭破齐。齐王烹郦生。东走高密。项羽闻韩信已举河北兵破齐。赵。且欲击楚。则使龙且。周兰往击之。韩信与战。骑将灌婴击。大破楚军。杀龙且。齐王广奔彭越。当此时。彭越将兵居梁地。往来苦楚兵。绝其粮食。\\
\hspace*{2em}四年。项羽乃谓海春侯大司马曹咎曰。﹁谨守成皋。若汉挑战。慎勿与战。无令得东而已。我十五日必定梁地。复从将军。﹂乃行击陈留。外黄。睢阳。下之。汉果数挑楚军。楚军不出。使人辱之五六日。大司马怒。度兵汜水。士卒半渡。汉击之。大破楚军。尽得楚国金玉货赂。大司马咎。长史欣皆自刭汜水上。项羽至睢阳。闻海春侯破。乃引兵还。汉军方围锺离眜於荥阳东。项羽至。尽走险阻。\\
\hspace*{2em}韩信已破齐。使人言曰。﹁齐边楚。权轻。不为假王。恐不能安齐。﹂汉王欲攻之。留侯曰。﹁不如因而立之。使自为守。﹂乃遣张良操印绶立韩信为齐王。\\
\hspace*{2em}项羽闻龙且军破。则恐。使盱台人武涉往说韩信。韩信不听。\\
\hspace*{2em}楚汉久相持未决。丁壮苦军旅。老弱罢转饟。汉王项羽相与临广武之间而语。项羽欲与汉王独身挑战。汉王数项羽曰。﹁始与项羽俱受命怀王。曰先入定关中者王之。项羽负约。王我於蜀汉。罪一。秦项羽矫杀卿子冠军而自尊。罪二。项羽已救赵。当还报。而擅劫诸侯兵入关。罪三。怀王约入秦无暴掠。项羽烧秦宫室。掘始皇帝冢。私收其财物。罪四。又彊杀秦降王子婴。罪五。诈阬秦子弟新安二十万。王其将。罪六。项羽皆王诸将善地。而徙逐故主。令臣下争叛逆。罪七。项羽出逐义帝彭城。自都之。夺韩王地。并王梁楚。多自予。罪八。项羽使人阴弑义帝江南。罪九。夫为人臣而弑其主。杀已降。为政不平。主约不信。天下所不容。大逆无道。罪十也。吾以义兵从诸侯诛残贼。使刑馀罪人击杀项羽。何苦乃与公挑战。﹂项羽大怒。伏弩射中汉王。汉王伤匈。乃扪足曰。﹁虏中吾指。﹂汉王病创卧。张良彊请汉王起行劳军。以安士卒。毋令楚乘胜於汉。汉王出行军。病甚。因驰入成皋。\\
\hspace*{2em}病愈。西入关。至栎阳。存问父老。置酒。枭故塞王欣头栎阳市。留四日。复如军。军广武。关中兵益出。\\
\hspace*{2em}当此时。彭越将兵居梁地。往来苦楚兵。绝其粮食。田横往从之。项羽数击彭越等。齐王信又进击楚。项羽恐。乃与汉王约。中分天下。割鸿沟而西者为汉。鸿沟而东者为楚。项王归汉王父母妻子。军中皆呼万岁。乃归而别去。\\
\hspace*{2em}项羽解而东归。汉王欲引而西归。用留侯。陈平计。乃进兵追项羽。至阳夏南止军。与齐王信。建成侯彭越期会而击楚军。至固陵。不会。楚击汉军。大破之。汉王复入壁。深堑而守之。用张良计。於是韩信。彭越皆往。及刘贾入楚地。围寿春。汉王败固陵。乃使使者召大司马周殷举九江兵而迎武王。行屠城父。随刘贾。齐梁诸侯皆大会垓下。立武王布为淮南王。\\
\hspace*{2em}五年。高祖与诸侯兵共击楚军。与项羽决胜垓下。淮阴侯将三十万自当之。孔将军居左。费将军居右。皇帝在後。绛侯。柴将军在皇帝後。项羽之卒可十万。淮阴先合。不利。卻。孔将军。费将军纵。楚兵不利。淮阴侯复乘之。大败垓下。项羽卒闻汉军之楚歌。以为汉尽得楚地。项羽乃败而走。是以兵大败。使骑将灌婴追杀项羽东城。斩首八万。遂略定楚地。鲁为楚坚守不下。汉王引诸侯兵北。示鲁父老项羽头。鲁乃降。遂以鲁公号葬项羽穀城。还至定陶。驰入齐王壁。夺其军。\\
\hspace*{2em}正月。诸侯及将相相与共请尊汉王为皇帝。汉王曰。﹁吾闻帝贤者有也。空言虚语。非所守也。吾不敢当帝位。﹂群臣皆曰。﹁大王起微细。诛暴逆。平定四海。有功者辄裂地而封为王侯。大王不尊号。皆疑不信。臣等以死守之。﹂汉王三让。不得已。曰。﹁诸君必以为便。便国家。﹂甲午。乃即皇帝位氾水之阳。\\
\hspace*{2em}皇帝曰义帝无後。齐王韩信习楚风俗。徙为楚王。都下邳。立建成侯彭越为梁王。都定陶。故韩王信为韩王。都阳翟。徙衡山王吴芮为长沙王。都临湘。番君之将梅鋗有功。从入武关。故德番君。淮南王布。燕王臧荼。赵王敖皆如故。\\
\hspace*{2em}天下大定。高祖都雒阳。诸侯皆臣属。故临江王驩为项羽叛汉。令卢绾。刘贾围之。不下。数月而降。杀之雒阳。\\
\hspace*{2em}五月。兵皆罢归家。诸侯子在关中者复之十二岁。其归者复之六岁。食之一岁。\\
\hspace*{2em}高祖置酒雒阳南宫。高祖曰。﹁列侯诸将无敢隐朕。皆言其情。吾所以有天下者何。项氏之所以失天下者何。﹂高起。王陵对曰。﹁陛下慢而侮人。项羽仁而爱人。然陛下使人攻城略地。所降下者因以予之。与天下同利也。项羽妒贤嫉能。有功者害之。贤者疑之。战胜而不予人功。得地而不予人利。此所以失天下也。﹂高祖曰。﹁公知其一。未知其二。夫运筹策帷帐之中。决胜於千里之外。吾不如子房。镇国家。抚百姓。给餽饟。不绝粮道。吾不如萧何。连百万之军。战必胜。攻必取。吾不如韩信。此三者。皆人杰也。吾能用之。此吾所以取天下也。项羽有一范增而不能用。此其所以为我擒也。﹂\\
\hspace*{2em}高祖欲长都雒阳。齐人刘敬说。乃留侯劝上入都关中。高祖是日驾。入都关中。六月。大赦天下。\\
\hspace*{2em}十月。燕王臧荼反。攻下代地。高祖自将击之。得燕王臧荼。即立太尉卢绾为燕王。使丞相哙将兵攻代。\\
\hspace*{2em}其秋。利几反。高祖自将兵击之。利几走。利几者。项氏之将。项氏败。利几为陈公。不随项羽。亡降高祖。高祖侯之颍川。高祖至雒阳。举通侯籍召之。而利几恐。故反。\\
\hspace*{2em}六年。高祖五日一朝太公。如家人父子礼。太公家令说太公曰。﹁天无二日。土无二王。今高祖虽子。人主也。太公虽父。人臣也。柰何令人主拜人臣。如此。则威重不行。﹂後高祖朝。太公拥篲。迎门卻行。高祖大惊。下扶太公。太公曰。﹁帝。人主也。柰何以我乱天下法。﹂於是高祖乃尊太公为太上皇。心善家令言。赐金五百斤。\\
\hspace*{2em}十二月。人有上变事告楚王信谋反。上问左右。左右争欲击之。用陈平计。乃伪游云梦。会诸侯於陈。楚王信迎。即因执之。是日。大赦天下。田肯贺。因说高祖曰。﹁陛下得韩信。又治秦中。秦。形胜之国。带河山之险。县隔千里。持戟百万。秦得百二焉。地埶便利。其以下兵於诸侯。譬犹居高屋之上建瓴水也。夫齐。东有琅邪。即墨之饶。南有泰山之固。西有浊河之限。北有勃海之利。地方二千里。持戟百万。县隔千里之外。齐得十二焉。故此东西秦也。非亲子弟。莫可使王齐矣。﹂高祖曰。﹁善。﹂赐黄金五百斤。\\
\hspace*{2em}後十馀日。封韩信为淮阴侯。分其地为二国。高祖曰将军刘贾数有功。以为荆王。王淮东。弟交为楚王。王淮西。子肥为齐王。王七十馀城。民能齐言者皆属齐。乃论功。与诸列侯剖符行封。徙韩王信太原。\\
\hspace*{2em}七年。匈奴攻韩王信马邑。信因与谋反太原。白土曼丘臣。王黄立故赵将赵利为王以反。高祖自往击之。会天寒。士卒堕指者什二三。遂至平城。匈奴围我平城。七日而後罢去。令樊哙止定代地。立兄刘仲为代王。\\
\hspace*{2em}二月。高祖自平城过赵。雒阳。至长安。长乐宫成。丞相已下徙治长安。\\
\hspace*{2em}八年。高祖东击韩王信馀反寇於东垣。\\
\hspace*{2em}萧丞相营作未央宫。立东阙。北阙。前殿。武库。太仓。高祖还。见宫阙壮甚。怒。谓萧何曰。﹁天下匈匈苦战数岁。成败未可知。是何治宫室过度也。﹂萧何曰。﹁天下方未定。故可因遂就宫室。且夫天子四海为家。非壮丽无以重威。且无令後世有以加也。﹂高祖乃说。\\
\hspace*{2em}高祖之东垣。过柏人。赵相贯高等谋弑高祖。高祖心动。因不留。代王刘仲弃国亡。自归雒阳。废以为合阳侯。\\
\hspace*{2em}九年。赵相贯高等事发觉。夷三族。废赵王敖为宣平侯。是岁。徙贵族楚昭。屈。景。怀。齐田氏关中。\\
\hspace*{2em}未央宫成。高祖大朝诸侯群臣。置酒未央前殿。高祖奉玉卮。起为太上皇寿。曰。﹁始大人常以臣无赖。不能治产业。不如仲力。今某之业所就孰与仲多。﹂殿上群臣皆呼万岁。大笑为乐。\\
\hspace*{2em}十年十月。淮南王黥布。梁王彭越。燕王卢绾。荆王刘贾。楚王刘交。齐王刘肥。长沙王吴芮皆来朝长乐宫。春夏无事。\\
\hspace*{2em}七月。太上皇崩栎阳宫。楚王。梁王皆来送葬。赦栎阳囚。更命郦邑曰新丰。\\
\hspace*{2em}八月。赵相国陈豨反代地。上曰。﹁豨尝为吾使。甚有信。代地吾所急也。故封豨为列侯。以相国守代。今乃与王黄等劫掠代地。代地吏民非有罪也。其赦代吏民。﹂九月。上自东往击之。至邯郸。上喜曰。﹁豨不南据邯郸而阻漳水。吾知其无能为也。﹂闻豨将皆故贾人也。上曰。﹁吾知所以与之。﹂乃多以金啗豨将。豨将多降者。\\
\hspace*{2em}十一年。高祖在邯郸诛豨等未毕。豨将侯敞将万馀人游行。王黄军曲逆。张春渡河击聊城。汉使将军郭蒙与齐将击。大破之。太尉周勃道太原入。定代地。至马邑。马邑不下。即攻残之。\\
\hspace*{2em}豨将赵利守东垣。高祖攻之。不下。月馀。卒骂高祖。高祖怒。城降。令出骂者斩之。不骂者原之。於是乃分赵山北。立子恆以为代王。都晋阳。\\
\hspace*{2em}春。淮阴侯韩信谋反关中。夷三族。\\
\hspace*{2em}夏。梁王彭越谋反。废迁蜀。复欲反。遂夷三族。立子恢为梁王。子友为淮阳王。\\
\hspace*{2em}秋七月。淮南王黥布反。东并荆王刘贾地。北渡淮。楚王交走入薛。高祖自往击之。立子长为淮南王。\\
\hspace*{2em}十二年。十月。高祖已击布军会甀。布走。令别将追之。\\
\hspace*{2em}高祖还归。过沛。留。置酒沛宫。悉召故人父老子弟纵酒。发沛中兒得百二十人。教之歌。酒酣。高祖击筑。自为歌诗曰。﹁大风起兮云飞扬。威加海内兮归故乡。安得猛士兮守四方。﹂令兒皆和习之。高祖乃起舞。慷慨伤怀。泣数行下。谓沛父兄曰。﹁游子悲故乡。吾虽都关中。万岁後吾魂魄犹乐思沛。且朕自沛公以诛暴逆。遂有天下。其以沛为朕汤沐邑。复其民。世世无有所与。﹂沛父兄诸母故人日乐饮极驩。道旧故为笑乐。十馀日。高祖欲去。沛父兄固请留高祖。高祖曰。﹁吾人众多。父兄不能给。﹂乃去。沛中空县皆之邑西献。高祖复留止。张饮三日。沛父兄皆顿首曰。﹁沛幸得复。丰未复。唯陛下哀怜之。﹂高祖曰。﹁丰吾所生长。极不忘耳。吾特为其以雍齿故反我为魏。﹂沛父兄固请。乃并复丰。比沛。於是拜沛侯刘濞为吴王。\\
\hspace*{2em}汉将别击布军洮水南北。皆大破之。追得斩布鄱阳。\\
\hspace*{2em}樊哙别将兵定代。斩陈豨当城。\\
\hspace*{2em}十一月。高祖自布军至长安。十二月。高祖曰。﹁秦始皇帝。楚隐王陈涉。魏安釐王。齐缗王。赵悼襄王皆绝无後。予守冢各十家。秦皇帝二十家。魏公子无忌五家。﹂赦代地吏民为陈豨。赵利所劫掠者。皆赦之。陈豨降将言豨反时。燕王卢绾使人之豨所。与阴谋。上使辟阳侯迎绾。绾称病。辟阳侯归。具言绾反有端矣。二月。使樊哙。周勃将兵击燕王绾。赦燕吏民与反者。立皇子建为燕王。\\
\hspace*{2em}高祖击布时。为流矢所中。行道病。病甚。吕后迎良医。医入见。高祖问医。医曰。﹁病可治。﹂於是高祖嫚骂之曰。﹁吾以布衣提三尺剑取天下。此非天命乎。命乃在天。虽扁鹊何益。﹂遂不使治病。赐金五十斤罢之。已而吕后问。﹁陛下百岁後。萧相国即死。令谁代之。﹂上曰。﹁曹参可。﹂问其次。上曰。﹁王陵可。然陵少戆。陈平可以助之。陈平智有馀。然难以独任。周勃重厚少文。然安刘氏者必勃也。可令为太尉。﹂吕后复问其次。上曰。﹁此後亦非而所知也。﹂\\
\hspace*{2em}卢绾与数千骑居塞下候伺。幸上病愈自入谢。\\
\hspace*{2em}四月甲辰。高祖崩长乐宫。四日不发丧。吕后与审食其谋曰。﹁诸将与帝为编户民。今北面为臣。此常怏怏。今乃事少主。非尽族是。天下不安。﹂人或闻之。语郦将军。郦将军往见审食其。曰。﹁吾闻帝已崩。四日不发丧。欲诛诸将。诚如此。天下危矣。陈平。灌婴将十万守荥阳。樊哙。周勃将二十万定燕。代。此闻帝崩。诸将皆诛。必连兵还乡以攻关中。大臣内叛。诸侯外反。亡可翘足而待也。﹂审食其入言之。乃以丁未发丧。大赦天下。\\
\hspace*{2em}卢绾闻高祖崩。遂亡入匈奴。\\
\hspace*{2em}丙寅。葬。己巳。立太子。至太上皇庙。群臣皆曰。﹁高祖起微细。拨乱世反之正。平定天下。为汉太祖。功最高。﹂上尊号为高皇帝。太子袭号为皇帝。孝惠帝也。令郡国诸侯各立高祖庙。以岁时祠。\\
\hspace*{2em}及孝惠五年。思高祖之悲乐沛。以沛宫为高祖原庙。高祖所教歌兒百二十人。皆令为吹乐。後有缺。辄补之。\\
\hspace*{2em}高帝八男。长庶齐悼惠王肥。次孝惠。吕后子。次戚夫人子赵隐王如意。次代王恆。已立为孝文帝。薄太后子。次梁王恢。吕太后时徙为赵共王。次淮阳王友。吕太后时徙为赵幽王。次淮南厉王长。次燕王建。\\
\hspace*{2em}太史公曰。夏之政忠。忠之敝。小人以野。故殷人承之以敬。敬之敝。小人以鬼。故周人承之以文。文之敝。小人以僿。故救僿莫若以忠。三王之道若循环。终而复始。周秦之间。可谓文敝矣。秦政不改。反酷刑法。岂不缪乎。故汉兴。承敝易变。使人不倦。得天统矣。朝以十月。车服黄屋左纛。葬长陵。\\
% ===== 索隐述赞（司马贞《史记索隐》）· 统一样式 =====
% 题头「索隐述赞曰」五字并入反白框；正文每四句一列，定宽盒 \hfil 居列中；
% \fontsize{24}{41} 列距对齐全书界行（121+41k）；\clearpage 使述赞独占一页。
\clearpage
\shuzanhead{索隐述赞曰}\linebreak
{\shuzanbody\fontsize{24}{41}\selectfont
\hbox to \hsize{\hfil 高祖初起。始自徒中。言从泗上。即号沛公。\hfil}\linebreak
\hbox to \hsize{\hfil 啸命豪杰。奋发材雄。彤云郁砀。素灵告丰。\hfil}\linebreak
\hbox to \hsize{\hfil 龙变星聚。蛇分径空。项氏主命。负约弃功。\hfil}\linebreak
\hbox to \hsize{\hfil 王我巴蜀。实愤于衷。三秦既北。五兵遂东。\hfil}\linebreak
\hbox to \hsize{\hfil 氾水即位。咸阳筑宫。威加四海。还歌大风。\hfil}%
\par}
\end{document}
